\chapter{Descripción del proyecto}
\section{Problema y solución}
En una plataforma cinematográfica, la gestión de cuentas, películas, comunidades, reacciones y sesiones compartidas responde a necesidades distintas y cambia a ritmos diferentes. Cuando estas funciones se concentran en una sola aplicación, las reglas de autenticación, la información del catálogo, la actividad social y la coordinación en tiempo real quedan estrechamente acopladas. Esto dificulta mantener cada parte, incorporar cambios y elegir la forma de persistencia más conveniente para sus datos.

Para abordar esta situación, ASTRA CLOUD adopta una arquitectura de microservicios. Identity se encarga del registro y la autenticación; Catalog administra la información de las películas; Community organiza perfiles, clubes y salas; Interaction registra reseñas, likes, listas e historial; Cinema Session coordina las sesiones temporales; y Analytics consulta la información reunida para fines analíticos. La separación permite delimitar responsabilidades y establecer integraciones explícitas entre los servicios, sin perder la continuidad de la experiencia para el usuario.
\section{Alcance}
El proyecto incluye las seis APIs de backend, sus configuraciones de contenedor, los procesos de carga hacia S3 y la plantilla de infraestructura. La mayor parte de las operaciones se resuelve mediante REST; Cinema Session complementa ese modelo con WebSocket para comunicar eventos de reproducción, participación y chat.
\section{Capa cliente}
La aplicación web está desplegada en AWS Amplify. Desde ella se puede iniciar sesión, explorar el catálogo, revisar el detalle de una película, consultar clubes, editar el perfil y participar en órbitas compartidas. La interfaz consume las APIs publicadas por API Gateway y archivos estáticos de S3. El código fuente del frontend y la configuración de Amplify no forman parte de este repositorio; no obstante, la aplicación desplegada se describe en el Capítulo~\ref{chap:frontend}.
